\section{Nontrivial affine-space fibrations}
\label{sec:affine-fibrations}

The residual-coordinate property of $H$ determines the fibers of two
natural morphisms. We first use it to prove that the morphisms are smooth,
then use the nonpolynomiality of $A$ to rule out Zariski local triviality.

\begin{corollary}[Nontrivial smooth $\Aff^3$-fibrations]
\label{cor:affine-fibrations}
Let $P,H,A$ be as in \eqref{eq:example}, and put $X=\Spec A$. Write
$\Aff^1_p=\Spec\kk[p]$ and
$\Aff^2_{p,q}=\Spec\kk[p,q]$, where $q$ is independent of $p$.
Consider the morphisms
\[
 \pi:X\longrightarrow\Aff^1_p,
 \qquad
 \rho=(p,H):\Spec P=\Aff^5\longrightarrow\Aff^2_{p,q},
\]
induced respectively by $\kk[p]\to A$ and by
$\kk[p,q]\to P$ with $q\mapsto H$.
Both are smooth surjective morphisms of relative dimension three.
For either morphism and every point $y$ of its base, the fiber over $y$
is isomorphic to $\Aff^3_{\kappa(y)}$; in particular, every geometric
fiber is affine three-space. Neither morphism is a
Zariski-locally trivial $\Aff^3$-bundle.
Thus the Dolgachev--Weisfeiler affine-fibration conjecture fails in
relative dimension three over each of the smooth complex affine bases
$\Aff^1$ and $\Aff^2$.
\end{corollary}

\begin{proof}
Set $\Lambda=\kk[p]$ and write $\mathbf z=(s,u,F,J)$. For every prime
$\mathfrak p$ of $\Lambda$, let $H_{\mathfrak p}$ denote the image of $H$ in
$\kappa(\mathfrak p)[\mathbf z]$. The coordinate changes in
\Cref{con:localized-coordinates,con:localized-h} show that $H$ is a
coordinate over $\kk[p,p^{-1}]$. At $\mathfrak p=(p)$, the determinant-one
change in \Cref{con:linear-coordinates,con:linear-inverse} gives
$H_{\mathfrak p}=L\bmod p$. Every other prime of $\Lambda$ avoids $p$.
It follows that
for every $\mathfrak p$ there are three polynomials $G_1,G_2,G_3$ over
$\kappa(\mathfrak p)$ such that
\begin{equation}\label{fib:residual-coordinate}
 \kappa(\mathfrak p)[\mathbf z]
 =\kappa(\mathfrak p)[H_{\mathfrak p},G_1,G_2,G_3]
\end{equation}
is a polynomial coordinate system.

Let $y$ be any point of $\Aff^2_{p,q}$, set
$\mathfrak p=\ker(\Lambda\to\kappa(y))$, and let $b$ be the image of $q$ in
$\kappa(y)$. The field $\kappa(y)$ is an extension of
$\kappa(\mathfrak p)$. Tensoring \eqref{fib:residual-coordinate} with
$\kappa(y)$ and then setting the coordinate $H_{\mathfrak p}$ equal to
$b$ gives
\[
 \rho_y
 =\Spec\bigl(\kappa(y)[\mathbf z]/(H_{\mathfrak p}-b)\bigr)
 \simeq\Aff^3_{\kappa(y)}.
\]
Here $H_{\mathfrak p}$ in the quotient denotes its image after this
field extension. Thus every fiber of $\rho$ is affine three-space over
its residue field, and the same is true after extension to an algebraic
closure. In particular, $\rho$ is surjective. Moreover,
\[
 A=P/(H)\simeq P\otimes_{\Lambda[q]}\Lambda,
 \qquad q\longmapsto0,
\]
so $\pi$ is the base change of $\rho$ along the closed immersion
$\Aff^1_p\hookrightarrow\Aff^2_{p,q}$ defined by $q=0$. The fiber
assertions and surjectivity follow for $\pi$ as well.

We next check smoothness without assuming local triviality. Let
\[
 \mathfrak j=
 \left(\frac{\partial H}{\partial s},
       \frac{\partial H}{\partial u},
       \frac{\partial H}{\partial F},
       \frac{\partial H}{\partial J}\right)\subset P.
\]
For every $\mathfrak p$, differentiation with respect to the coordinate
$H_{\mathfrak p}$ in \eqref{fib:residual-coordinate} gives a
$\kappa(\mathfrak p)$-derivation $\delta$ with
$\delta(H_{\mathfrak p})=1$. The chain rule expresses $1$ as a linear
combination of the four partial derivatives of $H_{\mathfrak p}$.
Consequently the image of $\mathfrak j$ generates the unit ideal in
$\kappa(\mathfrak p)[\mathbf z]$. If a prime of $P$ contained
$\mathfrak j$, localization and reduction over its contraction to $\Lambda$
would give a proper prime of this fiber algebra containing that unit
ideal. Hence $\mathfrak j=P$.

As a $\Lambda[q]$-algebra, $P$ has the presentation
\[
 P\simeq \Lambda[q][\mathbf z]/(q-H).
\]
The principal open sets on which one of the four displayed partial
derivatives is invertible cover $\Spec P$. On each such open set this
presentation has one relation with an invertible Jacobian minor, so it
is standard smooth of relative dimension $4-1=3$. Thus $\rho$ is smooth
of relative dimension three. Smoothness and relative dimension are
preserved by base change, giving the same assertion for $\pi$.

Suppose now that one of the two morphisms is Zariski-locally trivial,
and put
\[
 (R_0,B)=(\kk[p],A)
 \quad\hbox{or}\quad
 (R_0,B)=(\kk[p,q],P),
\]
respectively, with the algebra structures specified in the statement.
In either case, $B$ is a finitely presented $R_0$-algebra, and local triviality
implies
$B_{\mathfrak m}\simeq R_{0,\mathfrak m}\poly3$
for every maximal ideal $\mathfrak m$ of $R_0$.
Bass--Connell--Wright's theorem
\cite[Theorem~3(a)]{BassConnellWright}, in the formulation recalled by
Gupta~\cite[Theorem~4.1]{GuptaSurvey}, therefore gives an
$R_0$-algebra isomorphism
\[
 B\simeq\operatorname{Sym}_{R_0}(M)
\]
for a projective $R_0$-module $M$. This conclusion allows arbitrary
polynomial transition automorphisms in the assumed local trivializations;
no linear structure is assumed beforehand.

The module $M$ is finitely generated. Indeed, if $\mathfrak a$ is the
augmentation ideal of $\operatorname{Sym}_{R_0}(M)$, then
$M\simeq\mathfrak a/\mathfrak a^2$; finite generation follows because
$B$ is a finite-type augmented $R_0$-algebra. Tensoring with each residue
field and comparing the dimensions of the resulting polynomial rings
shows that $M$ has rank three. For the freeness step, Suslin's
Theorem~2(b)~\cite{Suslin} states that a finitely generated projective
module of rank greater than $\dim C$ over $C[t_1,\ldots,t_m]$ is extended
from $C$ when $C$ is regular and noetherian. Here $C=\kk$ has dimension
zero, giving the needed Quillen--Suslin statement: $M$ is extended from
a three-dimensional complex vector space.
Thus $M$ is free and
\[
 B\simeq R_0\poly3.
\]
For $\pi$ this would give
$A\simeq\kk[p]\poly3\simeq\kk\poly4$, contrary to \Cref{thm:main}.
For $\rho$ an $R_0$-algebra isomorphism would give a polynomial coordinate
system
$P=\kk[p,H,G_1,G_2,G_3]$ for suitable $G_1,G_2,G_3\in P$.
In particular, $H$ would be a coordinate of $P$, contrary to
Corollary~\ref{cor:stable-coordinate}. Neither morphism is therefore
Zariski-locally trivial.
\end{proof}

The two bases are smooth, normal, and affine over $\kk$. The fiber
calculation verifies the residue-field hypothesis in the original
affine-fibration assertion of Dolgachev--Weisfeiler, not only its
geometric-fiber version. These counterexamples have relative dimension
three; they do not address the corresponding assertions in relative
dimensions one or two.
